\chapter{Experiment 1: Exploratory Spectral Tensor Benchmark}
\label{ch:v1}

\section{Objective and input cohort}

The first experiment tests whether a compact non-negative
subject-channel-frequency-time tensor can support three-class prediction on the
controlled 69-person cohort. It uses the 12,270-window clean primary partition
before the later EEGLAB boundary correction. The experiment is valuable as a
baseline and design diagnostic, but its maximum across ranks and models is
exploratory because the outer score is also used to compare those choices.

\section{Tensor construction}

For every four-second window and each of 19 channels, Welch PSD uses a 500-sample
Hann segment, 250-sample overlap, 500-point FFT, constant detrending, and density
scaling. The 59 bins from 1 to 30 Hz are retained. Power is normalized over the
whole channel-frequency map and square-root transformed to obtain a non-negative
representation.

Each participant's windows are sorted chronologically and assigned to 10
rank-based time bins. The median spectral map within each bin produces

\begin{equation}
\mathcal{X}\in\mathbb{R}_{+}^{69\times19\times59\times10}.
\end{equation}

The 10 bins summarize early-to-late recording position; they are not event-
aligned epochs across people.

\section{Decomposition methods}

The benchmark evaluates ranks 3, 5, and 8 for five representations:

\begin{enumerate}
  \item PCA on the flattened tensor, fitted on training participants;
  \item NMF with NNDSVDa initialization, coordinate descent, 1,000 iterations,
  tolerance $10^{-5}$, and held-out transform;
  \item non-negative CP fitted by multiplicative updates;
  \item graph-regularized non-negative CP with nominal graph weight 0.1; and
  \item non-negative Tucker with subject rank equal to the experimental rank and
  feature ranks capped at channel 8, frequency 12, and time 5.
\end{enumerate}

CP/Tucker held-out features are solved in a fixed training-derived basis using
non-negative least squares. Decomposition fitting occurs separately inside each
subject-level training fold.

\section{Classifiers and validation}

Each representation is paired with regularized multinomial logistic regression,
RBF SVM, and histogram gradient boosting. Logistic and SVM use standardized
features and balanced class weights. The evaluation uses three repetitions of
stratified five-fold participant CV with seed 5040. Since each tensor row is one
participant, there is no window-level split at this stage.

\section{Results}

\begin{table}[htbp]
\centering
\begin{tabular}{lrrl}
\toprule
\textbf{Representation} & \textbf{Best rank} & \textbf{Mean BA} & \textbf{Associated classifier}\\
\midrule
PCA & 8 & 54.1\% & Logistic regression\\
NMF & 8 & 54.1\% & Logistic regression\\
Non-negative Tucker & 3 & 53.6\% & Histogram gradient boosting\\
Non-negative CP & 5/8 & 52.2\% & Logistic regression\\
Graph non-negative CP & 5 & 51.7\% & Logistic regression\\
\bottomrule
\end{tabular}
\caption{Strongest observed V1 result for each representation family. These are exploratory maxima, not nested-selection estimates.}
\label{tab:v1-results}
\end{table}

For the cited PCA rank-8 result, mean \BA{} is 54.1\% with repeat SD 3.6
percentage points. Its probability-averaged 69-subject confusion matrix has
recalls of 52.2\% \AD{}, 30.4\% \FTD{}, and 73.9\% \CN{}. The low \FTD{}
recall immediately shows that equal window counts do not eliminate the
difficulty of the smallest diagnostic group.

An NMF convergence check raises the rank-8 iteration limit to 5,000 and produces
the same classification result. Reconstruction error improves with rank for CP,
but prediction does not improve monotonically. This separates compression
fidelity from discriminative utility.

\section{Feature-bottleneck diagnostic}

A later diagnostic constructs less-compressed channel-band-time features and
uses repeated five-fold CV with training-only univariate selection. Its strongest
combination, 760 raw features reduced to 40 and classified by a linear SVM,
reaches 57.0\% mean balanced accuracy. This suggests that aggressive low-rank
compression may remove some useful structure, but it is not a causal ablation
because feature form and selection also change.

\section{Known limitations and correction path}

The first tensorizer attached a hard-coded electrode name order to rows that were
actually stored in another EEGLAB order. Consequently, the graph edges and any
electrode-specific interpretation in V1 are invalid. PCA, flattened NMF, global
statistics, and channel-permutation-invariant accuracy calculations are not
invalidated by a consistent row permutation. The graph-NCP objective also has a
factor-scale inconsistency. V2 therefore reads headers directly, uses signed
absolute logPSD, guards boundaries, and evaluates a fully nested pipeline.

